%% 05-related-work.tex --- included by ../sfc.tex
\section{Related Work}
\label{sec:related}

\textbf{\TeleSABRE{}}~\cite{russo2025telesabre} is the published state of
the art for this cost model, supporting teledata and gate teleportation
on \SABRE{}'s extended-set scoring. The two differ chiefly in decision
structure. \TeleSABRE{} scores inter-core candidates through a
contracted-communication-graph Dijkstra energy folding staging cost and
lookahead into one value, and --- uniquely here --- pools those
candidates with local SWAPs, taking the global minimum each iteration.
\PORTER{} keeps local work strictly ahead of inter-core work, scores
whole plans independently of preparation cost, generates preparation
SWAPs only for the plan it commits to, adds the packet primitive, and
treats capacity as a legality condition with its own relief search
rather than as \TeleSABRE{}'s safety-valve cap.

\textbf{\dSABRE{}}~\cite{lidsabre} attacks the same problem and is a close
relative. We run it in its default configuration, and our numbers match
those it publishes for the architectures it reports.

Stating what the two share keeps either from being credited with the
other's ideas. Both rank moves by \SABRE{}'s front-plus-lookahead
$\Delta$ score; both are strictly locality-first, scoring an inter-core
candidate only once no front gate is intra-core; both promote capacity
from a scoring penalty to a legality condition with the same floor of two
free slots; and both manufacture capacity by relaying a free slot from a
reservoir along a BFS path rather than evicting arbitrarily. Four
differences remain, and they are what this paper evaluates: \PORTER{}'s
action vocabulary is larger; its unit of commitment is a whole plan
rather than one hop; preparation cost enters as a tiebreaker rather than
additively; and its initial layout is one deterministic construction
rather than a best-of-$N$ search.

The fifth difference is not a design choice but a gap, and it runs
\dSABRE{}'s way. Three guarantees are worth separating.
\emph{Correctness}: the returned stream realises every gate and violates
no coupling or port rule. \emph{Termination}: every run ends, in success
or by raising a diagnostic or tripping a guard. \emph{Completion}: that
it ends in success. \PORTER{} provides the first two by construction
(Section~\ref{sec:sfc-progress}) and the third only empirically.
\dSABRE{}'s safe mode provides all three, carrying a
termination-with-success proof under stated architectural hypotheses ---
a connected core graph, connected intra-core couplings, and enough spare
capacity per core. \PORTER{} has no analogue; what it has is measured
rather than proved, and no instance in this evaluation is left
unrouted.

\begin{table}[t]
\centering
\caption{Design axes of the compilers discussed here. ``Commitment'' is
  the unit a scoring decision is taken over. \texttt{pytket-dqc} is
  partition-first and does not route with these primitives, so several
  axes do not apply.}
\label{tab:taxonomy}
\setlength{\tabcolsep}{3pt}
\small
\resizebox{\ifdim\width>\columnwidth \columnwidth\else\width\fi}{!}{%
\begin{tabular}{@{}l l l l l@{}}
\toprule
 & \TeleSABRE{} & \dSABRE{} & pytket-dqc & \PORTER{} \\
\midrule
Paradigm       & online & online & partition-first & online \\
Primitives     & data, gate & data & (e-bit) & data, gate, packet \\
Local/remote   & one pool & staged & --- & staged \\
Commitment     & one move & one hop & whole plan & whole plan \\
Prep.\ cost    & in score & in score & --- & tie-break \\
Layout         & Hungarian, & SabreLayout, & hypergraph & one \\
               & \ 3 seeds  & \ 3 seeds    & partition  & \ deterministic \\
Capacity       & cap & legality + relay & partition & legality + relief \\
Termination    & --- & with-success proof & --- & correct, terminates \\
Metric         & EPR & EPR & e-bits & EPR \\
\bottomrule
\end{tabular}%
}
\end{table}

Table~\ref{tab:taxonomy} places the four systems on the axes that
matter here.

\textbf{Other approaches.} Partition-first compilers assign qubits to
cores by hypergraph or community
partitioning~\cite{escofet2023hqa,pytketdqc2024,daei2020optimized} and
route the residue; \texttt{pytket-dqc}~\cite{pytketdqc2024} is the mature
representative, compared against in \appPytket{} under the caveat that
its e-bit unit assumes long-lived entanglement and different
communication memory. AutoComm~\cite{wu2022autocomm} identifies gate
sequences that can share one entangled resource, which \PORTER{}'s packet
applies inside the routing loop rather than as a pre-pass
(Section~\ref{sec:primitives}). \SABRE{}~\cite{li2019tackling} supplies
the scoring skeleton all three inherit.
